\Askhsh{24769}{4}
\begin{ylh}{1}
    \item 2.3 Κανόνες παραγώγισης
    \item 2.6 Συνέπειες του Θεωρήματος Μέσης Τιμής
    \item 2.7 Τοπικά ακρότατα συνάρτησης
    \item 2.8 Κυρτότητα - Σημεία καμπής συνάρτησης
    \item 3.1 Αόριστο ολοκλήρωμα
\end{ylh}
Δίνεται η συνάρτηση $f(x)=\ln(x+1)-\frac{x}{x+1}$, $x>-1$ και έστω F αρχική της f με $F(1)=\ln2$.
\begin{alist}
    \item Να αποδείξετε ότι για κάθε $x>-1$ ισχύει $f'(x)=\frac{x}{(x+1)^2}$ και να μελετήσετε τη συνάρτηση f ως προς τη μονοτονία. \monades{8}
    \item Να αποδείξετε ότι η F είναι κυρτή στο διάστημα $[0, +\infty)$. \monades{6}
    \item \begin{rlist}
        \item Να βρείτε την εξίσωση της εφαπτομένης της γραφικής παράστασης της F στο $x_0=1$. \monades{6}
        \item Να αποδείξετε ότι για κάθε $x>0$ ισχύει $\frac{2F(x)-1}{x} \ge \ln4-1$. \monades{5}
    \end{rlist}
\end{alist}